% Tapis de Quillen, with Quillen's offprint and preprint (cote 162-5, typed part): archive cote 162-5, site page 51 (PDF page 52).
% offprint. Draft compared: none.
% AI-transcribed by Claude Opus 5.5 (claude-opus-5-5, adaptive thinking, effort high), Message Batches API, batch t162_full (msgbatch_01UAh3zyWPd454AeDkhnmQ33), 2026-11-13; not yet checked by eye.
\scanpage{162-5/51}
\mspage{238} \hfill D.G. Quillen:

degree $k$ and $S_k T$ has the basis $z^\alpha$, $|\alpha|=k$ where $\alpha$ is a multi-index. If we
identify $V$ with $\mathbf{C}^r$ and let $e_1,\dots,e_r$ be the canonical basis of $\mathbf{C}^r$, then
$S_k T\otimes V$ can be naturally identified with the holomorphic polynomial
functions homogeneous of degree $k$ with values in $\mathbf{C}^r$. Then $S_m T\otimes V$
has the basis $z^\alpha e_i$, $|\alpha|=m$, $1\leqq i\leqq r$, so there are complex numbers $a_{\alpha i,\beta j}$,
$|\alpha|=|\beta|=m$, $1\leqq i, j\leqq r$ with $a_{\beta j,\alpha i}=a^*_{\alpha i,\beta j}$ such that
\[
\begin{aligned}
E(u_1,u_2)&=\sum \alpha!\,\beta!\,a_{\alpha i,\beta j}\,b_{\alpha i}\,c^*_{\beta j}\\
E^\wedge(\zeta,v)&=\sum a_{\alpha i,\beta j}\,\zeta^\alpha\,v_i\,\zeta^{*\beta}\,v^*_j
\end{aligned}
\tag{2}
\]
if $u_1=\sum b_{\alpha i}\,z^\alpha e_i$, $u_2=\sum c_{\alpha i}\,z^\alpha e_i$, $\zeta=\sum \zeta^i z^i$, and $v=(v_1,\dots,v_r)$, where *
denotes complex conjugate. The matrix $a_{\alpha i,\beta j}$ may be determined from
the Taylor series of $E^\wedge$ at $(0,0)$ which shows that $E$ and $E^\wedge$ determine
each other.

Let $E^\sim$ be the differential quadratic form
\[
E^\sim(f,g)=\sum a_{\alpha i,\beta j}(D^\alpha f_i)(D^\beta g_j)^*
\]
where $f$ and $g$ are holomorphic functions on $\mathbf{C}^n$ with values in $\mathbf{C}^r$ and
$D^i=\partial/\partial z^i$. If $dG=\pi^{-n}\exp(-|z|^2)\,dx^1\,dy^1\dots dx^n\,dy^n$ is the Gaussian
measure on $\mathbf{C}^n$, then we have the following expression for $E^N$.

Lemma 1. If $f,g\in S_{m+N}T\otimes V$, then
\[
E^N(f,g)=\int E^\sim(f,g)\,dG.
\]

Proof. If $u_1$, $u_2$ are two holomorphic polynomial functions on $\mathbf{C}^n$, let
\[
(u_1,u_2)_0=\int u_1\,u_2^*\,dG.
\]

By integrating by parts we obtain the basic formula
\[
(D^i u_1,u_2)_0=(u_1,z^i u_2)_0,
\]
from which follows
\[
(z^\alpha,z^\beta)_0=\begin{cases}0 & \alpha\neq\beta\\ \alpha! & \alpha=\beta.\end{cases}
\tag{3}
\]
By definition
\[
\begin{aligned}
(E^N)^\wedge(\zeta,v)&=(N!)^{-1}\,|\zeta|^{2N}\,E(\zeta,v)=\sum_{|\gamma|=N}(\gamma!)^{-1}\,\zeta^\gamma\,\zeta^{*\gamma}\,E(\zeta,v)\\
&=\sum(\gamma!)^{-1}\,a_{\alpha i,\beta j}\,\zeta^{\gamma+\alpha}\,v_i(\zeta^*)^{\gamma+\beta}\,v_j^*,
\end{aligned}
\]
hence
\[
(E^N)^\sim(f,g)=\sum_{|\gamma|=N}(\gamma!)^{-1}\,E^\sim(D^\gamma f,D^\gamma g).
\]
By linearity it suffices to prove the lemma when $f=z^\delta e_i$ and $g=z^\varepsilon e_j$
where $|\delta|=|\varepsilon|=m+N$. But
\[
E^N(f,g)=(E^N)^\sim(f,g)(0)=\sum_{\substack{|\alpha|=|\beta|=m\\ |\gamma|=N}}(\gamma!)^{-1}\,a_{\alpha i,\beta j}(D^{\gamma+\alpha}z^\delta)(D^{\gamma+\beta}z^\varepsilon)^*
\]

\mspage{\hand{51}}
